\section{Introduction}\label{sec:intro}
Webshells are the standard persistence mechanism after a web-application compromise. GlareShell~\paperref{}
proposes to detect PHP webshells statically by turning each script into an interprocedural control-flow
graph (ICFG) of PHP statements, attributing each node with a word embedding of the normalised statement and a
risk level, and classifying the graph with a GNN. It reports an F1 of 0.9887 with GAT, above sequence- and
tree-based approaches.

This report reproduces GlareShell from its description -- neither code nor data were released -- and then asks
whether the reported performance measures what it claims to measure. Our contributions:
\begin{enumerate}
  \item \textbf{A working re-implementation} on top of the same static-analysis tool (LChecker), including four
        extractor problems that silently corrupt the graphs if left unfixed (\Cref{sec:method}).
  \item \textbf{A rebuilt, documented corpus} (GlareShell-R) from the paper's sources with provenance, two
        deduplication levels, near-duplicate clusters, a data card and a label guide (\Cref{sec:dataset}).
  \item \textbf{A faithful replication} of the paper's Tables 2 and 4 and Figures 6--8 with five seeds and confidence
        intervals (\Cref{sec:repro}).
  \item \textbf{A leakage and generalisation analysis}: a cluster-disjoint split, a cross-source test on four unseen
        benign projects, simple baselines, and an error taxonomy (\Cref{sec:beyond,sec:errors}).
\end{enumerate}

\paragraph{Summary of findings}
\begin{itemize}
  \item \textbf{RQ1:} the headline performance is reproduced: F1 0.98 on the paper's protocol (paper: 0.97--0.99).
  \item \textbf{RQ2:} the component claims are not: GAT $\approx$ GCN $\approx$ GIN, embedding dimension and context
        window barely matter, and normalisation/risk weights do not significantly help (paired tests, $p \geq 0.07$).
  \item \textbf{RQ3:} 77\% of test webshells leak via near-duplicates (F1 $-$2 points without them); a token
        TF-IDF + LR baseline is better than every GNN, and graph statistics alone reach F1 0.92.
  \item \textbf{RQ4:} on four unseen benign projects the FPR rises from 0.4\% to 1.7\%; recall on webshells from
        new collections is 0.54--0.64, and only 6\% at a 0.1\% alarm budget; 10 of 12 small legitimate scripts are flagged. The detector has learned
        ``not framework code'' rather than ``malicious''.
  \item \textbf{Improvement:} training on five small benign applications lowers the GAT FPR from 2.5\% to 0.9\%
        on all benign test files ($p=0.006$, recall unchanged), but 8 of the 12 scripts remain flagged.
\end{itemize}

